% Abstrak berbahasa Inggris. Isinya harus setara dengan abstrak Indonesia.
\begin{abstract}
\panduan{The abstract must be a single paragraph (200--250 words, without citations). Encapsulate empirical problem background, objective, methodology, metric results, and conclusion. Minimum 5 keywords. Disable this box with panduan=false.}

Usability evaluation of public service information systems is essential to ensure effective technology adoption and citizen satisfaction within village government administrations. The implementation of village digital services frequently encounters obstacles due to citizens' diverse digital literacy and suboptimal system interfaces. This study aims to evaluate the usability of a web-based village demographic administration service application using the System Usability Scale (SUS) method and formulate empirical evidence-based interface improvement recommendations. The research adopts a quantitative approach involving 30 village community respondents selected via purposive sampling. Empirical data were gathered using the standardized 10-item Likert-scale SUS questionnaire (scales 1--5), validated for reliability, combined with structured task observations and semi-structured user interviews. The evaluation yielded an overall average SUS score of 74.50 with a standard deviation of 8.25. Benchmarked against standard SUS interpretation frameworks, this score places the application in the Acceptable range, achieves a Good adjective rating, and attains a Grade B percentile ranking, confirming operational viability. Item-level analysis revealed that functional integration ease received the highest ratings, whereas requirement document upload workflows and initial technical guidance needs received the lowest evaluations. Consequently, interface recommendations focus on streamlining digital form layouts, incorporating contextual interactive tooltips, and enforcing real-time client-side input validation. The study concludes that the application is usable and well received by citizens, with the projected implementation of recommendations reducing village administrative transaction duration by up to 30%.

\keywords{usability evaluation; information system; system usability scale; village administration; user satisfaction}
\end{abstract}
